\subsection{Tipos de nodos y contratos de ejecución}
\label{sec:nodos-contratos}

En LangGraph todo paso del pipeline es un nodo: una función Python con la firma
\texttt{(state: AgentState) $\to$ Command}.
El nodo recibe el estado compartido completo, ejecuta su lógica —ya sea una llamada a un LLM,
un cálculo determinista o una suspensión HITL— y devuelve un \texttt{Command} que especifica
dos cosas: qué campos del estado deben actualizarse (\texttt{update=\{\ldots\}}) y cuál es el nodo
destino (\texttt{goto="workflow\_controller"} en todos los casos salvo el controlador).
LangGraph aplica la actualización parcial al estado antes de invocar el nodo siguiente;
ningún nodo puede leer su propia escritura en la misma ejecución.

La tabla~\ref{tab:nodos-resumen} resume el tipo y la responsabilidad de cada nodo.

\begin{figure}[H]
\centering
\begin{tabular}{lll}
\hline
\textbf{Nodo} & \textbf{Tipo} & \textbf{Responsabilidad principal} \\
\hline
\texttt{workflow\_controller}   & Determinista (máquina de estados)  & Enrutamiento centralizado \\
\texttt{data\_validator}        & Agente ReAct                        & Validación y limpieza del dataset \\
\texttt{dataset\_approval}      & Puerta HITL (gate 1)                & Aprobación humana del dataset \\
\texttt{planner}                & Agente ReAct + salida estructurada  & Planificación de modelos candidatos \\
\texttt{executor}               & Determinista (Optuna + MLflow)      & Entrenamiento y selección del campeón \\
\texttt{evaluation}             & Determinista                        & Comparación con umbrales y campeón actual \\
\texttt{report\_writer}         & LLM salida estructurada             & Informe de auditoría \\
\texttt{deployment\_approval}   & Puerta HITL (gate 2)                & Aprobación humana del despliegue \\
\texttt{deployer}               & Determinista (MLflow)               & Registro y promoción a campeón \\
\hline
\end{tabular}
\caption{Resumen de nodos del grafo MLOps.}
\label{tab:nodos-resumen}
\end{figure}

A continuación se describe cada nodo con sus campos de entrada (lectura del estado),
sus campos de salida (escritura mediante \texttt{Command.update}) y los detalles de su implementación.

% ---------------------------------------------------------------------------
\subsubsection*{workflow\_controller}
\label{sec:nodo-controller}

\textbf{Tipo:} determinista --- máquina de estados Python sin llamadas a LLM.\\
\textbf{Fichero:} \texttt{src/mlops\_agents/graphs/workflow\_controller.py}

\textbf{Entradas del estado leídas:}
\begin{itemize}
  \item \texttt{error\_message} — si está poblado, redirige a \texttt{END} inmediatamente.
  \item \texttt{validation\_passed} — \texttt{False} implica invocar \texttt{data\_validator}.
  \item \texttt{agent\_attempt\_counts} — conteo de intentos por nodo; bloquea si supera el máximo configurable.
  \item \texttt{dataset\_approved} — \texttt{None} implica pasar a la puerta 1; \texttt{False} implica reintentar \texttt{data\_validator}.
  \item \texttt{training\_plan} — ausente implica invocar \texttt{planner}.
  \item \texttt{training\_run\_id} — ausente implica invocar \texttt{executor}.
  \item \texttt{evaluation\_passed} — \texttt{None} implica invocar \texttt{evaluation}; \texttt{False} termina el pipeline.
  \item \texttt{evaluation\_report\_audit} — ausente implica invocar \texttt{report\_writer}.
  \item \texttt{deployment\_approved} — \texttt{None} implica pasar a la puerta 2; \texttt{False} termina el pipeline.
  \item \texttt{deployment\_decision} — \texttt{"pending"} implica invocar \texttt{deployer}.
\end{itemize}

\textbf{Salidas (campos escritos):}
Solo en el caso de rechazo por intentos máximos escribe \texttt{error\_message}.
En todos los demás casos el controlador no escribe ningún campo; únicamente decide el \texttt{goto}.

El controlador es el único nodo del grafo que lee la totalidad del estado compartido.
Todos los demás nodos acceden únicamente a los campos que necesitan para su función específica.

% ---------------------------------------------------------------------------
\subsubsection*{data\_validator}
\label{sec:nodo-data-validator}

\textbf{Tipo:} agente ReAct (\texttt{create\_agent()}) con modelo \texttt{gpt-5.4-mini}.\\
\textbf{Fichero:} \texttt{src/mlops\_agents/graphs/mlops\_graph.py} (\texttt{data\_validator\_node})

\textbf{Entradas del estado leídas:}
\begin{itemize}
  \item \texttt{dataset\_paths} — lista de rutas a los ficheros CSV brutos subidos por el usuario.
  \item \texttt{schema\_json} — esquema JSON con \texttt{problem\_type}, \texttt{target\_column} y metadatos de tarea.
  \item \texttt{dataset\_rejection\_comment} — comentario del operador humano en caso de reintento tras rechazo en la puerta 1.
\end{itemize}

\textbf{Comportamiento:} Antes de invocar al agente, una función determinista (\texttt{\_validate\_schema\_contract}) verifica el contrato del esquema. Si el esquema no cumple los requisitos mínimos (tipo de problema declarado, columna objetivo existente, campos de predicción temporal cuando el tipo lo requiere), el nodo escribe \texttt{error\_message} y retorna sin invocar al LLM. Si el contrato se cumple, el agente recibe el contexto de entrada y ejecuta su bucle de herramientas sobre el conjunto \{\texttt{load\_dataset}, \texttt{merge\_datasets}, \texttt{apply\_column\_mapping}, \texttt{validate\_against\_schema}, \texttt{check\_missing\_values}, \texttt{check\_data\_quality}, \texttt{impute\_missing\_values}, \texttt{parse\_datetime\_column}, \texttt{detect\_temporal\_gaps}\}. El nodo extrae el resultado de las herramientas \texttt{check\_data\_quality}, \texttt{apply\_column\_mapping}, \texttt{validate\_against\_schema} e \texttt{impute\_missing\_values} del historial de mensajes y construye los campos de salida.

\textbf{Salidas (campos escritos):}
\begin{itemize}
  \item \texttt{validation\_passed} — \texttt{True} si \texttt{validate\_against\_schema} devolvió \texttt{passed: true}.
  \item \texttt{validation\_report} — JSON de \texttt{check\_data\_quality} (calidad del dataset: valores ausentes, duplicados, etc.).
  \item \texttt{processed\_dataset\_path} — ruta al CSV canónico tras limpieza e imputación.
  \item \texttt{dataset\_summary} — \texttt{\{row\_count, column\_names, dtypes, null\_counts\}} calculado sobre el CSV canónico.
  \item \texttt{problem\_type} — extraído del esquema (\texttt{"classification"}, \texttt{"regression"} o \texttt{"forecasting"}).
  \item \texttt{task\_metadata} — \texttt{target\_column} más metadatos de predicción temporal si aplica.
  \item \texttt{schema\_json} — propagado al estado para auditoría.
  \item \texttt{dataset\_rejection\_comment} — vaciado a \texttt{""} tras una validación exitosa.
  \item \texttt{error\_message} — escrito solo si la validación falla tras el intento de autocorrección del agente.
\end{itemize}

% ---------------------------------------------------------------------------
\subsubsection*{dataset\_approval}
\label{sec:nodo-gate1}

\textbf{Tipo:} puerta HITL (gate 1) --- suspensión del grafo mediante \texttt{interrupt()}.\\
\textbf{Fichero:} \texttt{src/mlops\_agents/graphs/approval\_nodes.py}

\textbf{Entradas del estado leídas:}
\begin{itemize}
  \item \texttt{processed\_dataset\_path} — para construir la vista previa del dataset (5 primeras y últimas filas).
  \item \texttt{validation\_report} — informe de calidad incluido en el payload de la interrupción.
  \item \texttt{agent\_attempt\_counts} — número de intento actual, expuesto al operador.
  \item \texttt{problem\_type} — determina si se incluyen las filas del final (solo en predicción temporal).
\end{itemize}

\textbf{Comportamiento:} Construye una vista previa del dataset canónico y llama a \texttt{interrupt(\{\ldots\})}. La ejecución del grafo queda suspendida y el estado persiste en el \emph{checkpointer}. Cuando el operador envía su decisión vía \texttt{POST /runs/\{run\_id\}/approve}, LangGraph reanuda el nodo con el valor de \texttt{Command(resume=\{approved, comment\})}.

\textbf{Salidas (campos escritos):}
\begin{itemize}
  \item \texttt{dataset\_approved} — \texttt{True} o \texttt{False} según la decisión del operador.
  \item \texttt{dataset\_rejection\_comment} — comentario del operador si rechazó, vacío si aprobó.
\end{itemize}

% ---------------------------------------------------------------------------
\subsubsection*{planner}
\label{sec:nodo-planner}

\textbf{Tipo:} agente ReAct con salida estructurada (\texttt{create\_agent(\ldots, response\_format=PlannerOutput)}) con modelo \texttt{gpt-5.4-mini}.\\
\textbf{Fichero:} \texttt{src/mlops\_agents/planning/node.py}

\textbf{Entradas del estado leídas:}
\begin{itemize}
  \item \texttt{processed\_dataset\_path} — para construir el perfil del dataset (\texttt{DatasetProfile}).
  \item \texttt{problem\_type} — tipo de problema; determina qué modelos del registro son válidos.
  \item \texttt{task\_metadata} — metadatos de la tarea (columna objetivo, horizonte, frecuencia).
\end{itemize}

\textbf{Comportamiento:} El nodo construye determinísticamente un \texttt{PlannerValidationContext} (universo de modelos válidos, experiencias similares pre-recuperadas, reglas aplicables) antes de lanzar el agente. El agente dispone de cuatro herramientas: \texttt{list\_available\_models}, \texttt{retrieve\_similar\_experiences}, \texttt{retrieve\_ml\_knowledge} e \texttt{inspect\_model\_details} (máximo tres llamadas). Tras el bucle de herramientas, el agente produce un \texttt{PlannerOutput} validado por cuatro comprobaciones deterministas: integridad del plan (sin solapamientos candidato/descartado), exhaustividad del registro (todos los modelos válidos están o seleccionados o justificadamente descartados), coherencia de evidencias (ninguna referencia citada inexistente en el contexto) y resolución de conflictos (si se detectaron conflictos, el plan debe incluir su resolución). Si alguna verificación falla, el nodo reintenta una vez con el mensaje de error como retroalimentación. Un envoltorio de gestión de errores (\texttt{\_planner\_node\_with\_error\_handling}) captura cualquier \texttt{PlannerError} o \texttt{OutputParserException} del segundo intento y escribe \texttt{error\_message} en lugar de propagar la excepción.

\textbf{Salidas (campos escritos):}
\begin{itemize}
  \item \texttt{training\_plan} — \texttt{TrainingPlan} serializado como diccionario Pydantic: candidatos con prioridad e hiperparámetros iniciales, modelos descartados con justificación, presupuesto de trials Optuna.
  \item \texttt{planner\_analysis} — texto generado por el planificador que explica la síntesis de evidencias.
  \item \texttt{planner\_evidence\_used} — lista de referencias a experiencias y reglas utilizadas.
  \item \texttt{planner\_warnings} — lista de advertencias emitidas por el planificador.
  \item \texttt{planner\_status} — \texttt{"ok"}, \texttt{"retry\_ok"} o \texttt{"failed"}.
  \item \texttt{planner\_retry\_used} — \texttt{True} si fue necesario el segundo intento.
  \item \texttt{planner\_tool\_trace} — traza serializada de las llamadas a herramientas realizadas por el agente.
  \item \texttt{planner\_validation\_context} — subconjunto auditado del contexto de validación determinista.
  \item \texttt{\_planner\_output\_record} — registro interno completo consumido por el ejecutor y la SSE.
  \item \texttt{error\_message} — escrito únicamente si el planificador falla tras agotar los reintentos.
\end{itemize}

% ---------------------------------------------------------------------------
\subsubsection*{executor}
\label{sec:nodo-executor}

\textbf{Tipo:} determinista --- Optuna para búsqueda de hiperparámetros, MLflow para registro de experimentos.\\
\textbf{Fichero:} \texttt{src/mlops\_agents/graphs/mlops\_graph.py} (\texttt{executor\_node})

\textbf{Entradas del estado leídas:}
\begin{itemize}
  \item \texttt{training\_plan} — \texttt{TrainingPlan} serializado: lista de candidatos, presupuesto de trials.
  \item \texttt{processed\_dataset\_path} — CSV canónico sobre el que entrenar.
  \item \texttt{problem\_type} y \texttt{task\_metadata} — tipo de problema y metadatos de la tarea.
  \item \texttt{\_planner\_output\_record} — salida del planificador, registrada junto al experimento MLflow.
\end{itemize}

\textbf{Comportamiento:} Deserializa el \texttt{TrainingPlan} con validación Pydantic y llama a \texttt{run\_training\_plan}, que divide el dataset, entrena cada modelo candidato con Optuna, registra cada trial en MLflow y selecciona el modelo con mejor métrica de validación como campeón. Al finalizar, serializa el registro de experiencia en un fichero JSON en disco.

\textbf{Salidas (campos escritos):}
\begin{itemize}
  \item \texttt{training\_run\_id} — identificador del experimento padre en MLflow.
  \item \texttt{training\_metrics} — métricas del modelo campeón sobre el conjunto de validación.
  \item \texttt{champion\_candidate} — diccionario con \texttt{model\_key} e hiperparámetros del modelo ganador.
  \item \texttt{trained\_model\_path} — ruta al artefacto del modelo campeón.
  \item \texttt{train\_pool\_path}, \texttt{test\_path}, \texttt{split\_metadata\_path} — rutas a las particiones generadas.
  \item \texttt{experience\_record\_path} — ruta al fichero JSON con el registro de experiencia completo.
\end{itemize}

% ---------------------------------------------------------------------------
\subsubsection*{evaluation}
\label{sec:nodo-evaluation}

\textbf{Tipo:} determinista --- comparación de métricas con umbrales fijos y campeón actual.\\
\textbf{Fichero:} \texttt{src/mlops\_agents/evaluation/promotion.py}

\textbf{Entradas del estado leídas:}
\begin{itemize}
  \item \texttt{problem\_type} — determina la métrica primaria (F1 en clasificación, RMSE en regresión y predicción temporal) y los umbrales de promoción.
  \item \texttt{training\_metrics} — métricas del candidato sobre el conjunto de validación.
  \item \texttt{training\_run\_id} — usado para consultar el campeón actual en MLflow.
\end{itemize}

\textbf{Comportamiento:} Recupera las métricas del campeón actual desde MLflow, aplica los umbrales definidos para el tipo de problema y determina si el candidato supera al campeón. La decisión es completamente determinista: no hay LLM en este nodo.

\textbf{Salidas (campos escritos):}
\begin{itemize}
  \item \texttt{evaluation\_passed} — \texttt{True} si el candidato supera los umbrales y al campeón actual, \texttt{False} en caso contrario.
  \item \texttt{candidate\_metrics} — métricas del candidato copiadas para auditoría.
  \item \texttt{champion\_metrics} — métricas del campeón actual recuperadas de MLflow.
  \item \texttt{thresholds\_applied} — umbrales utilizados para la decisión.
  \item \texttt{evaluation\_report} — diccionario consolidado con candidato, campeón y \texttt{candidate\_run\_id}.
\end{itemize}

% ---------------------------------------------------------------------------
\subsubsection*{report\_writer}
\label{sec:nodo-report-writer}

\textbf{Tipo:} LLM con salida estructurada (\texttt{with\_structured\_output(EvaluationReport)}) con modelo \texttt{gpt-5.4-nano}.\\
\textbf{Fichero:} \texttt{src/mlops\_agents/evaluation/report\_writer.py}

\textbf{Entradas del estado leídas:}
\begin{itemize}
  \item \texttt{problem\_type}, \texttt{task\_metadata} — contexto del problema.
  \item \texttt{training\_plan}, \texttt{training\_run\_id}, \texttt{training\_metrics} — qué se entrenó y con qué resultado.
  \item \texttt{candidate\_metrics}, \texttt{champion\_metrics}, \texttt{thresholds\_applied} — resultado de la evaluación determinista.
  \item \texttt{planner\_analysis}, \texttt{planner\_evidence\_used} — razonamiento previo del planificador.
  \item \texttt{evaluation\_report\_audit} (ausente en este punto) — vacío; el nodo lo produce.
\end{itemize}

\textbf{Comportamiento:} Invoca el LLM con un prompt de sistema fijo y un mensaje de contexto ensamblado determinísticamente a partir de los campos anteriores. El LLM produce un \texttt{EvaluationReport} en formato estructurado sin bucle de herramientas. Si el primer intento falla, el nodo reintenta una vez con el mensaje de error como retroalimentación. Si el segundo intento también falla, produce un informe de emergencia (\emph{stub}) para no bloquear el pipeline.

\textbf{Salidas (campos escritos):}
\begin{itemize}
  \item \texttt{evaluation\_report\_audit} — \texttt{EvaluationReport} serializado: análisis narrativo de consistencia entre planificación y resultados, veredicto, y observaciones sobre las experiencias previas.
  \item \texttt{evaluation\_report\_audit\_status} — \texttt{"ok"}, \texttt{"retry\_ok"} o \texttt{"stub"}.
\end{itemize}

% ---------------------------------------------------------------------------
\subsubsection*{deployment\_approval}
\label{sec:nodo-gate2}

\textbf{Tipo:} puerta HITL (gate 2) --- suspensión del grafo mediante \texttt{interrupt()}.\\
\textbf{Fichero:} \texttt{src/mlops\_agents/graphs/approval\_nodes.py}

\textbf{Entradas del estado leídas:}
\begin{itemize}
  \item \texttt{evaluation\_report}, \texttt{evaluation\_report\_audit} — informes determinista y narrativo para presentar al operador.
  \item \texttt{candidate\_metrics}, \texttt{champion\_metrics}, \texttt{thresholds\_applied} — métricas comparativas.
  \item \texttt{training\_plan}, \texttt{training\_run\_id} — identificación del candidato a desplegar.
\end{itemize}

\textbf{Comportamiento:} Construye el payload de la interrupción incluyendo todos los informes relevantes y el resumen de la acción propuesta (\texttt{register\_and\_promote} con alias \texttt{champion}). Llama a \texttt{interrupt(\{\ldots\})} y suspende el grafo hasta que el operador envíe su decisión.

\textbf{Salidas (campos escritos):}
\begin{itemize}
  \item \texttt{deployment\_approved} — \texttt{True} o \texttt{False} según la decisión del operador. Un valor \texttt{False} es terminal: el controlador redirige a \texttt{END} sin reintento.
\end{itemize}

% ---------------------------------------------------------------------------
\subsubsection*{deployer}
\label{sec:nodo-deployer}

\textbf{Tipo:} determinista --- registro y promoción en MLflow Model Registry.\\
\textbf{Fichero:} \texttt{src/mlops\_agents/deployment/deployer.py}

\textbf{Entradas del estado leídas:}
\begin{itemize}
  \item \texttt{training\_run\_id} — identificador del experimento MLflow con el modelo campeón a registrar.
\end{itemize}

\textbf{Comportamiento:} Invoca \texttt{register\_model} con el \texttt{run\_id} del campeón para crear una nueva versión en el Model Registry de MLflow, a continuación invoca \texttt{set\_model\_alias} para asignarle el alias \texttt{champion}. Ambas operaciones utilizan las herramientas MLflow del sistema.

\textbf{Salidas (campos escritos):}
\begin{itemize}
  \item \texttt{deployment\_status} — \texttt{"deployed"}.
  \item \texttt{deployment\_decision} — \texttt{"deployed"} (sobreescribe el valor previo \texttt{"pending"}).
  \item \texttt{best\_model\_uri} — URI canónica del modelo en el registro, con formato \texttt{models:/\{nombre\}/\{versión\}}.
\end{itemize}

Tras la escritura de estos campos, el controlador lee \texttt{deployment\_decision != "pending"} y redirige a \texttt{END}, completando la ejecución del pipeline.
